% Main mathematical alternative with an explicit suspended application.
% Requires amsmath, amssymb. No document wrapper or custom environments.

\section{A single suspended-application rule}

The cleanest formulation introduces an explicit suspended application
$U(a;u.h)$.  Its numerical meaning is $h[a/u]$, but its body is not
executed by term normalization.  Iteration at a syntactic successor
counter is factored into one such application.  The only arithmetic
inference then evaluates an application whose argument is known equal
to another term:
\[
 \boxed{\qquad
 \frac{\Gamma\vdash a=b\qquad\Gamma\vdash U(a;u.h)=r}
      {\Gamma\vdash r=N(h[b/u])}\quad\mathrm{App}.
 \qquad}
\]
Reflexivity is retained.  Unlike the baseline Step rule, App does not
inspect an iterator counter or reconstruct a predecessor iteration.
All of that syntactic work is in the specified canonical constructor.
The constructor, the added syntax, and their validation must be counted
in the complete proof checker; this section makes no size claim by
counting the inference rule alone.

\subsection{Syntax and normalization}

Raw terms are
\[
 t::=x\mid0\mid St\mid U(a;u.h)\mid I(n,a;u.h).
\]
In both binding constructors, $u$ is bound in $h$ only, not in $a$ or
$n$.  The application constructor is a definitional extension of the
old language:
\[
             U(a;u.h)^\circ:=I(S0,a^\circ;u.h^\circ).
\]
The map $(-)^\circ$ otherwise preserves variables, zero, successor,
and iteration.  In the ordinary scalar theory, one-step computation
gives $U(a;u.h)^\circ=h^\circ[a^\circ/u]$.  Thus the extra constructor
has no extra numerical freedom or axiomatic strength.

On already canonical arguments define
\begin{align*}
 K(0,a;u.h)&=a,\\
 K(Sn,a;u.h)&=U(K(n,a;u.h);u.h),\\
 K(n,a;u.h)&=I(n,a;u.h)
       &&\text{when $n$ has neither root $0$ nor root $S$.}
\end{align*}
The recursive call in the middle clause is on the proper subterm $n$
of $Sn$.  Define $N$ structurally by
\begin{align*}
 N(x)&=x,&N(0)&=0,&N(St)&=S N(t),\\
 N(U(a;u.h))&=U(N(a);u.N(h)),\\
 N(I(n,a;u.h))&=K(N(n),N(a);u.N(h)).
\end{align*}
The canonical terms are exactly the fixed points of $N$.  Equivalently,
all children are canonical and the counter of every retained iterator
has neither root $0$ nor root $S$.  A suspended application is always
retained, even when its body is constant or the identity.

For instance,
\[
 N(I(SSx,a;u.h))
 =U\bigl(U(I(x,N(a);u.N(h));u.N(h));u.N(h)\bigr).
\]
Each occurrence of $u$ binds only its own displayed body.  The result
is obtained by inspecting the leading successors of the counter,
without executing $h$ or computing any unknown numerical argument.

Extend $N$ through formulas, clauses, and contexts by normalizing their
numerical terms.  The formula constructors and their bound variables
are unchanged.  Canonical substitution is capture-free substitution
followed by $N$.

\subsection{A terminating and confluent rewrite specification}

The canonical forms are also the normal forms of two unconditional
contextual rewrite schemes:
\[
 I(0,a;u.h)\longrightarrow a,
 \qquad
 I(Sn,a;u.h)\longrightarrow U(I(n,a;u.h);u.h).
 \tag{ZF}
\]
The schemes are allowed under every binder, with the usual
capture-free convention.  There is deliberately no rewrite for $U$.
The two left-hand sides are left linear and disjoint at their roots;
all proper potential overlaps lie in the metavariables $n,a,h$.
The following direct argument supplies both termination and confluence
without relying on a particular higher-order rewriting formalism.

Assign positive integer weights by
\begin{align*}
 w(x)&=w(0)=1,&w(St)&=w(t)+1,\\
 w(U(a;u.h))&=w(a)+w(h)+1,\\
 w(I(n,a;u.h))&=(w(n)+1)(w(a)+w(h)+1).
\end{align*}
For the zero rewrite, the decrease is from
$2(w(a)+w(h)+1)$ to $w(a)$.  Put $A=w(a)$, $H=w(h)$, and $q=w(n)$.
For the successor rewrite, the left-hand side has weight
$(q+2)(A+H+1)$ and the right-hand side has weight
\[
             (q+1)(A+H+1)+H+1.
\]
The difference is exactly $A>0$.  Every constructor is strictly
increasing in each child's positive weight, so these decreases persist
in every context, including bodies.  Thus every reduction strictly
decreases a positive integer.  Duplication of the body causes no gap in
the argument: its cost is included explicitly in the calculation.

The function $N$ is invariant under either rewrite.  The zero case is
its defining clause; the successor case follows from
\[
 N(I(Sn,a;u.h))
 =U(K(N(n),N(a);u.N(h));u.N(h))
 =N(U(I(n,a;u.h);u.h)).
\]
Structural induction makes this invariance contextual.  Every term
reduces to an irreducible term by termination, and the irreducible terms
are exactly the fixed points of $N$.  Any such reduct of $t$ must
therefore equal $N(t)$.  Normal forms are unique, and any two reducts
of $t$ can be further reduced to this same normal form.  This proves
confluence as well as idempotence $N(N(t))=N(t)$.

For an ordinary primitive-recursive tree coding, the weight and one
deterministically chosen rewrite step are primitive recursive.  There
are fewer than $w(t)$ steps.  Bounded iteration of that step yields a
primitive-recursive implementation of $N$.  The direct smart-constructor
implementation is another such implementation.  This is a syntactic
normalizer, not a uniform evaluator for arbitrary primitive-recursive
programs: it never substitutes an argument into the body of $U$.

\subsection{Substitution and scoping}

Both rewrite schemes in (ZF) are preserved by every capture-free
substitution.  Rewriting the images of a substitution inside a fixed
term is also a contextual rewriting.  Confluence consequently gives
\[
 N(t[\sigma])=N\bigl(N(t)[N\circ\sigma]\bigr).
 \tag{SC}
\]
Under $U$ and the body of $I$, use the lifted substitution that fixes
their bound variable.  This is ordinary renaming and substitution
bookkeeping; factorization creates no new variable references.
In particular, $N$ commutes with renaming and introduces no new free
variables.  It need not preserve every free variable, since zero
contraction can discard an unused body.

Writing $t\langle\sigma\rangle=N(t[\sigma])$, (SC) and ordinary
capture-free substitution composition give
\[
 (t\langle\sigma\rangle)\langle\tau\rangle
 =t\langle x\mapsto\sigma(x)\langle\tau\rangle\rangle.
\]
The same statements hold for formulas and clauses, with their binders
lifted correctly.  They supply the exact coherence needed by the
logical inference rules and by the proof translation below.

\subsection{Equality and generic iterator computation from App}

Keep the axiom schema $\Gamma\vdash t=t$.  For
$J_s:=U(s;v.v)$, two applications of App give
\[
 \frac{s=t\qquad
       \displaystyle\frac{s=r\qquad J_s=J_s}{J_s=r}\ \mathrm{App}}
      {r=t}\ \mathrm{App}.
 \tag{Pivot}
\]
Hence symmetry and transitivity are available.  For any fixed numerical
template $h(v)$, use $H_s:=U(s;v.h(v))$.  App with $s=s$ and
$H_s=H_s$ gives $H_s=N(h[s/v])$; App with $s=t$ and $H_s=H_s$
gives $H_s=N(h[t/v])$.  Pivot proves congruence in $h$.  Combined
with another equality it gives the baseline crossed transport:
\[
 s=t,\qquad N(h[s/v])=r\quad\vdash\quad r=N(h[t/v]).
\]
No induction or numerical case distinction is used here.

For a canonical counter $c$, initializer $a$, and body $h$, the
canonical successor iteration is exactly
\[
                K(Sc,a;u.h)=U(K(c,a;u.h);u.h).
\]
Reflexivity on $K(c,a;u.h)$ and on this application, followed by App,
proves
\[
 K(Sc,a;u.h)=N(h[K(c,a;u.h)/u]).
\]
Thus arbitrary successor computation is recovered, despite the absence
of a counter in the arithmetic inference rule.  More generally, a
baseline Step instance with premises
$I(c,a;u.h)=t$ and $I(Sc,a;u.h)=r$ translates to a single App
instance, since the normalized second source is the application of
$h$ to the normalized first source.

\subsection{Logical rules and preservation of their modes}

The remaining logic is exactly the baseline five-scheme clause
interface: formulas are equations and
$D_{x,y}B=\exists x(B(x,0)\to\forall y\,B(x,y))$; contexts contain
reusable clauses $[A;y.B]$ with reading $A\to\forall yB$.  Together
with Refl and App, retain Use, DE, and the single witness-directed DI.
Use applies a retained clause to a proof of its antecedent.  DE opens
the witness of a $D$ formula as a fresh parameter and retains its
conditional-universal clause.  DI discharges the assumption $B(t,\eta)$
from a proof of $B(t,Sz)$, where $z$ is fresh and
\[
 \eta=z\quad\text{exactly when $B$ is a literal equation and $t=0$
 syntactically};\qquad \eta=0\text{ otherwise}.
\]
As always, branch contexts are weakened to the enlarged variable scope,
and a fresh branch variable cannot occur in its outer conclusion.
All substitutions use $N$.  The report must display the three full
inference schemes and all these scope conditions; they are summarized
here only to isolate the new arithmetic argument.

Let $N_0$ be the baseline's zero-only normalizer on its old language.
For an old-language term $t$,
\[
                   N(t)=0\quad\Longleftrightarrow\quad N_0(t)=0.
                   \tag{Zero}
\]
The proof is structural.  Variables and successors cannot become zero.
An iterator becomes zero exactly when its normalized counter becomes
zero and its normalized initializer becomes zero.  In every other
case its result is either iterator-headed or application-headed.
Use the induction hypotheses for the counter and initializer.

In particular, the new transformation maps an already zero-normal
witness to literal zero exactly when that witness was literal zero.
It also preserves formula root constructors.  Therefore DI uses the
same case or progress mode in a zero-normal baseline derivation and
in its translated derivation.  No new repair of the introduction mode
is needed.  This observation is stronger and more precise than
incorrectly assuming every syntax-dependent rule is preserved by every
arbitrary substitution.

\subsection{Both translations and both interpretation round trips}

Let $C_0$ be the baseline zero-normal clause calculus with Refl, Step,
Use, DE, DI, and let $C_U$ use the canonical terms above with Refl,
App, Use, DE, DI.  Translate old terms by $N$, formulas componentwise,
and contexts clausewise.  The preceding arguments give
\[
 C_0:\Gamma\vdash A
       \quad\Longrightarrow\quad
 C_U:N\Gamma\vdash NA.
\]
Refl is immediate; each Step is one App; Use and DE use (SC);
DI uses (SC) and (Zero).  Freshness is preserved because $N$ introduces
no variables.  This proof applies to arbitrary retained contexts.

In the other direction expand each $U$ by the map $(-)^\circ$ into
$I(1,a;u.h)$.  The ordinary scalar theory proves
\[
                      t^\circ=N(t)^\circ
\]
for every raw new-language term.  Each zero rewrite is the zero
computation equation.  Each successor factorization is proved by
comparing the two sides with $h[I(n,a;u.h)/u]$: the left uses successor
computation, and the right uses the one-step and zero equations.
Ordinary term congruence handles contexts except for pointwise changes
inside iterator bodies.  There one explicitly uses the baseline's
derived body-congruence lemma, proved by literal-equation induction on
the iteration counter.  This is an inherited dependency of the
normalization equivalence, even though App's equality bootstrap uses
no induction.

App is sound under this expansion by one-step computation and equality.
The logical rules retain their baseline ordinary meanings, because a
substitution followed by $N$ has provably the same numerical value as
the raw substitution.  The baseline ordinary/native equivalence
therefore supplies the reverse consequence simulation.  In
particular, a canonical sequent is provable precisely when its
expansion is provable in the baseline.

The syntax round trips are explicit.  If $t$ is canonical, then
\[
                         N(t^\circ)=t
\]
literally up to bound renaming.  Prove this structurally: an expanded
application $I(1,a^\circ;u.h^\circ)$ normalizes to $U(a;u.h)$,
while a retained iterator has a neutral counter whose normalized
expansion remains that same counter.  Conversely, an old term $s$
is provably equal in the baseline to $N(s)^\circ$.  This is the value
equality just proved.  Formula round trips follow by the ordinary
logical congruence rules and the inherited Eq/$D$ bridge.

Consequently the whole-domain bi-interpretation with full first-order
iPRA transfers.  The new constructor is a fixed primitive-recursive
abbreviation, and normalization changes expression representatives
only.  Numerical equality remains actual equality; no numerical
quotient, proper subdomain, or standard-model restriction is introduced.
Both composite numerical isomorphisms have graph $x=y$.  The
contradiction $0=S0$ is unchanged, so equiconsistency follows as well.

\subsection{What is and is not an improvement}

The mathematical change removes every iterator-counter condition from
the arithmetic inference rule.  All arithmetic reasoning uses one
application-evaluation scheme together with reflexivity; the full
calculus still has five schemes.  The two unconditional normalization
equations have a short termination proof, are stable under substitution,
and never evaluate a body.  This is simpler mathematically than keeping
$U$ as an abbreviation for a privileged unit iterator and excluding
that unit from the factorization rewrite.

The cost is one new term constructor and its canonical handling.  A
complete implementation can be larger or smaller than the current
Step checker; only a matched whole-module measurement settles that
question.  Neither deleting a counter from a displayed rule nor
moving computation into syntax establishes an absolute minimum.
